% CVT Foundations v0.2 — proposed Abstract and Introduction
% Review artifact only. This file does not replace main.tex.

\begin{abstract}
We propose Coherence Viability Theory (CVT) as a testable framework for boundary-mediated transformations in which the continued viability of an identifiable receiving or hosting system is part of the success criterion. CVT hypothesizes a candidate non-compensatory structure comprising bounded coupling, preserved identity-defining organization, pre-existing capture readiness, and available restorative reserve. These elements are not individually new: the framework draws from viability theory, autopoiesis, passivity-based control, resilience theory, and basin dynamics. Its proposed contribution is their conjunction as conditions of hosted transformation, together with a distinction between exchange delivered to a system and exchange retained as productive uptake. The non-compensability claim is conceptually prior to any particular scalar aggregator; the product-gate form developed here is one candidate representation rather than an established universal equation. We develop a minimal gated architecture and a membrane formulation, then use artificial fusion as an illustrative toy-model instantiation of the difference between initiating a transition and hosting it. CVT does not yet establish that the proposed conditions are universally necessary or sufficient. It instead defines a bounded research program in which the conditions, observables, aggregation rules, and counterexamples can be compared empirically or in reproducible simulation.
\end{abstract}

\section{Introduction}

Many systems depend on exchange across a boundary. They must receive energy, matter, information, influence, or control while retaining the organization that allows them to continue functioning as hosts. Neither closure nor openness alone solves this problem. Insufficient exchange can starve a system; excessive or poorly received exchange can overload it, erase relevant distinctions, exhaust recovery capacity, or drive it outside an admissible regime.

This paper studies a bounded class of such processes: \emph{boundary-mediated transformations in which continued host viability is part of what counts as success}. We call a transformation \emph{hosted} when a declared host undergoes or supports a meaningful transition while remaining viable over a stated observation horizon. The host, boundary, transformation variable, success criterion, timescale, and level of analysis must therefore be specified before CVT is applied.

CVT proposes that hosted transformation may depend on four jointly necessary conditions:
\begin{enumerate}[label=(\roman*)]
    \item \textbf{bounded coupling}: exchange is present but constrained relative to the host's changing handling range;
    \item \textbf{preserved distinction}: the host retains the identity-defining organization or independent constraints relevant to the declared success criterion;
    \item \textbf{capture readiness}: before uptake is observed, the host occupies a state from which delivered input can enter and remain within an admissible basin;
    \item \textbf{available restorative reserve}: sufficient internal or reliably accessible recovery capacity remains to absorb strain, repair damage, and preserve control authority.
\end{enumerate}

In compact form, the candidate structure is
\begin{equation}
    \mathrm{CVT}_{\mathrm{candidate}}
    =
    B\cap Q\cap C\cap S.
\end{equation}
This expression is a hypothesis about essential conditions, not a proof that the conditions are universal or sufficient. Their satisfaction defines a candidate viability region. Hidden state variables, delayed damage, environmental change, path dependence, adversarial interference, and measurement error may still produce failure.

A central distinction is between what is delivered and what is productively retained. Let \(J_{\mathrm{raw}}\) denote boundary flux delivered to the host and \(J_{\mathrm{eff}}\) denote retained productive uptake over a declared horizon. Realized uptake efficiency may be observed as
\begin{equation}
    \rho_{\mathrm{obs}}
    =
    \frac{J_{\mathrm{eff}}}{J_{\mathrm{raw}}},
\end{equation}
when \(J_{\mathrm{raw}}>0\). This ratio is an outcome, not an explanation. CVT must predict it from independently measured host conditions rather than define receptivity through the same gates later used to explain it. The testable claim is that increased delivered exchange can, in some regimes, reduce productive uptake when capture readiness or restorative reserve degrades sufficiently.

The individual ingredients of this framework are inherited from mature traditions. Viability theory studies admissible trajectories under state constraints. Autopoiesis studies self-producing organization and structural coupling. Passivity and port-Hamiltonian control formalize constrained energetic exchange. Resilience theory studies persistence and recovery under disturbance. Basin dynamics studies nonlinear capture and escape. CVT does not claim these concepts as new. Its proposed contribution is narrower: it asks whether a class of boundary-mediated failures is better understood when bounded coupling, identity preservation, capture readiness, and restorative reserve are treated as non-compensatory rather than substitutable.

Non-compensability is conceptually prior to the product equation used later in the manuscript. A product of normalized gates is one smooth representation because a near-zero essential gate sharply lowers the aggregate. It is not yet established as the uniquely correct form. Minimum-gate, weighted geometric-mean, normalized soft-minimum, and simultaneous constraint-intersection formulations remain competing representations to be compared against the same data or simulations.

The membrane formulation provides the paper's clearest structural illustration. A boundary can permit flux while the receiving domain remains unable to retain that flux productively. Permeability is therefore not identical to capture readiness or effective uptake. The fusion section extends this logic as an illustrative toy model: ignition access is distinguished from hosted burn, but the numerical thresholds and controller comparisons are not presented as calibrated reactor-physics results or as validation of CVT.

The paper proceeds as follows. Section~\ref{sec:lineage} situates CVT within its principal intellectual lineages. Section~\ref{sec:invariant} states the candidate four-condition structure and the non-compensability proposal. Section~\ref{sec:formal} develops an initial gated formal architecture whose variables and aggregation rules remain subject to operational testing. Section~\ref{sec:membrane} gives the membrane formulation. Section~\ref{sec:fusion} presents the fusion material as an illustrative instantiation of hosted transformation. The present paper establishes neither universal necessity nor sufficiency; it defines a research framework whose strongest claims must survive counterexamples, independent measurement, competing models, and delayed-failure tests.

